\begin{figure}[htbp]
\centering
\resizebox{0.92\textwidth}{!}{%
\begin{tikzpicture}[font=\scriptsize, node distance=8mm]
  \tikzset{evostep/.style={draw, rounded corners, align=center, minimum width=36mm, minimum height=10mm},
           version/.style={draw, circle, align=center, minimum size=13mm, fill=lightaccent}}

  \node[evostep, fill=warnbg] (known) {Ismert követelmények\\és kezdeti elképzelés};
  \node[evostep, fill=lightaccent, right=15mm of known] (initial) {Kezdeti\\implementáció};
  \node[evostep, fill=black!5, right=15mm of initial] (feedback) {Felhasználói\\véleményezés};
  \node[evostep, fill=lightaccent, right=15mm of feedback] (refine) {Finomítás, javítás,\\új funkciók};
  \node[evostep, fill=warnbg, right=15mm of refine] (target) {Megfelelő\\rendszer};

  \draw[arrow] (known) -- (initial);
  \draw[arrow] (initial) -- (feedback);
  \draw[arrow] (feedback) -- (refine);
  \draw[arrow] (refine) -- (target);
  \draw[arrow, dashed] ([yshift=-2mm]refine.south) .. controls +(0,-18mm) and +(0,-18mm) .. node[below, align=center]{új verzió, újabb visszajelzés} ([yshift=-2mm]initial.south);

  \node[version, below=13mm of initial] (v1) {v1};
  \node[version, below=13mm of feedback] (v2) {v2};
  \node[version, below=13mm of refine] (v3) {v3};
  \draw[arrow] (v1) -- (v2);
  \draw[arrow] (v2) -- (v3);

  \node[draw, rounded corners, fill=black!5, below=10mm of v2, align=center, minimum width=108mm] {Az evolúciós modellben a specifikáció, a fejlesztés és a validáció részben összemosódik: a rendszer a verziók során közelít a valódi igényekhez.};
\end{tikzpicture}%
}
\caption{Az evolúciós fejlesztés alapgondolata}
\label{fig:zv25_evolucios_modell}
\end{figure}
